\section{Setup}
\label{sec:setup}

We analyse a DreamerV3 world model \citep{hafner2023dreamerv3} trained on video of a pendulum. The
model is the reference PyTorch implementation at its published architecture: 13.5M parameters,
categorical $32\times32$ stochastic latents, a convolutional encoder and decoder, Kullback-Leibler (KL) balancing, and
unimix. We train the world model alone, with no actor and no critic, because we never need a policy.

\paragraph{Relation to our earlier study.} We developed this extraction procedure on a much
smaller system, and we describe it here in full because there is no separate paper to cite. A gated
recurrent unit (GRU) of about 13k parameters was trained to predict the next position of a simulated
oscillator from position alone, so it read the state directly instead of inferring it from pixels.
On that substrate the same search recovered a quantity correlating with true energy at $0.972$ or
better, scored close to zero on damped controls, and lowered rollout error by 68\% when enforced.
That study is also where the three controls in this paper come from: an untrained network, a
dissipative system, and a random law of matched complexity.

Two of its findings are retested in Section~\ref{sec:limits} and neither survives here. We treat
everything from it as a hypothesis to retest rather than as support, for two reasons. It is our own
unpublished work, so a reader cannot check it. And the two substrates do not share an error metric:
the GRU's rollout error lives in the oscillator's two-dimensional state space, while ours is pixel
error against held-out video. Percentages from the two are not interchangeable, and
Section~\ref{sec:intervention} says so where it quotes one.
Appendix~\ref{sec:gru} describes that study in full, including the results that did not survive it.

\paragraph{Code and run logs.} Every experiment in this paper, the run logs its figures are drawn
from, and the source of the paper itself are at
\url{https://github.com/Zarand3r/latent-noether-paper1}. The figures regenerate from the committed
logs with no GPU, so each number can be checked against the run that produced it. Model checkpoints
and rendered video are too large to commit; the repository pins the upstream commit of the DreamerV3
implementation our adapter is built against and gives the commands to rebuild both.

\paragraph{Data.} We render the video ourselves from gymnasium's \texttt{Pendulum-v1}
\citep{towers2024gymnasium}. Each step produces a $500\times500$ RGB frame, which we crop to
$448\times448$ and block-average to $64\times64$. We write initial conditions directly into the
environment state, sampling $\theta$ uniformly on $[-1.8, 1.8]$ and $\dot\theta$ on $[-2.2, 2.2]$,
and we reject any trajectory that reaches the simulator's $|\dot\theta| = 8$ speed clip, since that
clip is a hard nonlinearity that would destroy conservation. Actions are always zero, so the
pendulum evolves freely.

Velocity never appears in a single frame. There is no motion blur, no trail, and no torque arrow,
so the model has to infer $\dot\theta$ from the frame sequence on its own. That is what makes the
latent worth examining.

\paragraph{Training.} We train on a fixed dataset of 256 trajectories of 120 frames, using Adam at
$10^{-4}$, batch 16, sequence length 64, for a wall-clock budget of half an hour. That budget buys
about 6{,}500 gradient steps, a small fraction of a standard DreamerV3 run. We set free bits to 0
rather than the reference default of 1.0; the measured KL settles near 1.8 nats, far from collapse,
so the floor does no work here. Three seeds train independently and all three pass acceptance
checks we registered beforehand: KL above 1 nat, one-step decoding at least $4\times$ better than
predicting the dataset mean, and rollouts that stay finite. A model that failed to train would tell
us nothing about the analysis, so we discard it rather than report it.

This regime is offline, world-model-only, zero-action, and short. Our claims are about a DreamerV3
world model fitted to pendulum video, and not about DreamerV3 as a reinforcement learning agent.

\paragraph{Extraction.} After training we freeze the model. Everything that follows reads the
hidden state and never updates a weight. We encode the held-out videos, take the deterministic
recurrent state $h$, and project it onto its top 12 principal directions, dropping any direction the
data does not occupy. Inside that subspace we take the flow $F$ to be the model's own one-step
transition displacement, $T(h) - h$, rather than a vector field we fit ourselves. Fitting a field and
then analysing it would interpret our own regression instead of the model.

We then search for a degree-4 polynomial $C$ that is constant along the model's trajectories. The
search is a generalized eigenproblem: the invariance ratio, within-trajectory variance over total
variance, is a Rayleigh quotient, so one eigendecomposition returns the whole ranked family. We fit
$C$ jointly with an antisymmetric $B$ satisfying $F \approx B\nabla C$, which is the condition that
$C$ generates the flow rather than merely staying constant along it \citep{arnold1989mathematical}.
Section~\ref{sec:limits} reports how much that second condition contributes, which is less than we
expected.

\paragraph{What we compare against.} We score $C$ against the textbook energy
$E = \tfrac{1}{6}\dot\theta^2 + 5\cos\theta$. Gymnasium integrates semi-implicitly, so it conserves a
shadow Hamiltonian $\tilde H = H + O(\Delta t)$ rather than $E$ exactly
\citep{hairer2006geometric}. We measure about 12\% relative oscillation in $E$ from that mismatch,
with no secular drift. Ground truth enters only when we score a result, never inside the search.
